% Transcription — fonds Alexandre Grothendieck, shelfmark 156-6, batch 1 (pages 1-20).
% Established from the facsimile archives/batches/156-6/batch-01.pdf (a copy of
% 156-6.pdf fetched through the project's relay; no cover sheet in the batch,
% so PDF page k = folder page k), per the skill transcribe-grothendieck. The
% folder is ninety-three pages in five batches; this is the first. Batch 2
% (pages 21-40) and the first batch of the first draft (folder 156-4, GF IV)
% were read after drafting and the notation aligned with them.
%
% Pass: Opus 5.5 (claude-opus-5-5), 2026-09-26 — first pass, unchecked against the pages by a human.
%
% Pages skipped: none. Pages summarised: none. All twenty leaves are
% mathematics in his hand and all are transcribed.
%
% His pagination 1-20 runs across the top of the leaves and coincides with the
% archivists' (the « 19 » is written like a 17). Page 1 carries the heading
% « Analysis situs (deuxième mouture) » and, top right, « GF VI ». Dates in his
% hand, each given in a \note where it occurs: « 18 juin 86 », slanted in the
% left margin of page 1; « (20 juin) », heading an addition beside (1.16) b) on
% page 10; « le 20 » in the (mostly unread) margin of page 13; « 19 juin »,
% underlined, opening the last paragraph of page 20. The 20 June additions on
% pages 10 and 13 were made after the 19 June text of page 20. The bracketed
% « [Chapitre] » of the inventory title is the archivists'.
%
% Numbered runs: section « (I) » (circled) « Algèbre de figures » ou
% « Atelier » and paragraph « 1. » on page 1; his formula numbers (1.1)-(1.38)
% run through the batch (renumbered in places: (1.6) over (1.4) on p. 4, (1.7)
% over (1.5); two different formulas are both numbered (1.27) on p. 12; a
% struck (1.32) on p. 14), batch 2 continuing at (1.39); unnumbered starred
% maps (*), (**) on pp. 14, 19, 20. Conditions: « C 1, C 2, C 3 » of GF IV
% (§§ 5, 6, 8, read doubtfully) recalled on p. 16 with a margin renaming them
% « At 1, At 2, At 3 »; C 4 (margin) and C 5 on p. 19, a C 6 there struck out.
% Cross-references to the first draft (GF IV = folder 156-4): « Dans la
% première mouture » (p. 13) and « Dans GF IV (Analysis situs, première
% mouture) » (p. 16) — noted where they occur.
%
% What the batch is about. The second draft of Analysis situs restarts from an
% « algèbre de figures » or « atelier »: three sets, lieux ℒ ⊂ multistrates 𝓜
% ⊂ figures 𝔉 (1.1-1.2), with the déploiement F ↦ F̃ ⊂ 𝓜 (set of strata, the
% « incidence » relation) and F ↦ Fig_𝓜(F) (1.3-1.5), and the (grande) ombre
% Omb_𝓜(F) (the multistrates refining F) and multiombre Multomb (1.6-1.8);
% restricted to ℒ, the petite ombre omb and multomb (1.9), with the remark that
% a figure is recovered from its multiombre, not its ombre. Pages 7-12: six
% relations on 𝔉 — orders ≤ (sous-figure), ≪ (raffinement), ≼ (subdivision),
% ◁ (incidence, corrected in a margin to transitive non-reflexive), and the
% symmetric ∥ (disjonction) and compatibility (1.10-1.19), arranged in the
% implication diagram (1.18); which of them determine which (1.20-1.27): ≤ and
% ◁ each other, ≤ or ≪ each give ∥, ≪ gives ≼ and ℒ (the ≪-minimal non-empty
% figures), ≪ with compatibility gives ≤. Page 13 compares the three minimal
% presentations (1.28), recalls that the first draft used (≤, ≼) with ≪ as the
% generated relation (1.29), and settles on (≤, ≪) (1.30), with a 20 June
% margin « Non, ça ne suffit pas ». Pages 14-18: the same in terms of 𝓜 — 𝔉 as
% a préadmissible family of parts of 𝓜 with ≪ (1.33), conditions (i)-(ii), the
% C 1-C 3 of GF IV, finite-type figures (1.34), and the « moral » equivalence
% with (𝓜, ≤, compatibility) (1.35-1.36); aim: axioms valid both for
% Fig(L) of a set L and for finite affine simplicial figures over ℝ. Pages
% 19-20: Omb, Multomb in terms of 𝓜; C 4, C 5 (1.37) so that the multiombre is
% an ensemblist figure indexed by F̃; the injection 𝔉 → Fig(𝓜) and (1.38); on
% 19 June he drops the reading of G ≪ F as « sub-figure of a subdivision of F »
% for tame simplices, and the page breaks off mid-sentence.
%
% Notation fixed once, aligned with batch 2 and with chapters V (156-5) and IV
% (156-4): 𝔉 \mathfrak{F} (his script F), 𝓜 \mathcal{M}, ℒ \mathcal{L},
% 𝔓 \mathfrak{P} (once a roman P on p. 10, kept); his compatibility sign, an
% hourglass of two chevrons, \between (as batch 2); his ≤ with a curved lower
% stroke (subdivision) \preccurlyeq; incidence ◁ \triangleleft, its
% underlined form \trianglelefteq; his split lozenge \Diamond. « Omb »,
% « Multomb », « Dépl », « Fig » are his abbreviations, set in roman. His
% underlining is \emph in running text (\underline only inside math, p. 7).
% Circled letters and numerals are given in parentheses with a note.
%
% Hard pages: 13 and 19 (long slanted margin notes, read only in part), 15-16
% (last lines a tangle of strikes and interlinear insertions, order not
% certain), 20 (the ombre argument, interlinear). Variances on the page kept
% and flagged: « Dépl(F) » for Fig_𝓜(F) (p. 3); Multomb(F) = {X ∈ F̃ | Omb X}
% (p. 4, (1.7)) and « Multomb(X) » announced as Omb(X) (1.8); « raffinement »
% left after ≪ was replaced by ≼ (p. 9); « X ◁ 𝔉 » for X ◁ F (p. 11); « F = F »
% in the margin definition of the underlined triangle (p. 11); « par (1.25) »
% where ℒ is (1.26) (p. 14); « relation de disjonction » for the compatibility
% sign (p. 18).
\documentclass[11pt,a4paper]{article}
\input{../preamble/grothendieck}

\folder{156-6}
\batch{1}
\pages{1}{20}
\dating{1986}
\watermark{Édition de démonstration}
\foldertitle{[Chapitre] VI. Analysis situs (deuxième mouture) : notes manuscrites (18-20/06/1986)}

\begin{document}
\page{1}
\section{Analysis situs (deuxième mouture)}

\note{titre de sa main en tête de la page 1, qui porte en haut au centre son numéro « 1 » ; dans l'angle supérieur droit, dans un coin tracé à la plume, « GF VI » (le chiffre souligné)}

\marginal{18 juin 86}
\note{date de sa main, écrite en oblique dans la marge gauche, au début du premier alinéa}

Avant de décrire ce qu'est une \emph{contrée}, je vais décrire ce qui sera une « contrée avec notion de multistrates » — la famille des multistrates \add{choisies} jouant un peu le rôle de la famille d'ouverts \struck{baignant} \add{engendrant} \struck{\uncertain{pour}} une topologie, ou une famille génératrice d'éléments d'un topos. Je vais donc commencer par

\subsection{(I) « Algèbre de figures » ou « \emph{Atelier} ».}
\note{le « I » est cerclé}

1. Une \emph{algèbre de figures} implique avant tout trois types d'objets, les \emph{lieux}, les \emph{multistrates}, les \emph{figures}, formant trois ensembles
\[
  (1.1.) \qquad \mathcal{L},\ \mathcal{M},\ \mathfrak{F} ,
\]
liés entre eux par diverses applications, et chacun muni de diverses structures. Ainsi on a des applications canoniques injectives
\[
  (1.2) \qquad \mathcal{L} \overset{b)}{\hookrightarrow} \mathcal{M} \overset{a)}{\hookrightarrow} \mathfrak{F} ,
\]
que nous utiliserons souvent pour identifier un lieu à une multistrate particulière, et une multistrate à une figure particulière — $\mathcal{L}$ à un sous-ens.\ de $\mathcal{M}$, $\mathcal{M}$ à un sous-ens.\ de $\mathfrak{F}$.
\note{les étiquettes « b) » et « a) » sont écrites au-dessus des deux flèches, dans cet ordre, de gauche à droite}

Il y a d'autre part deux autres paires d'applications, que voici.

\page{2}
\note{son numéro « 2 » en haut de la page}
\[
  (1.3) \qquad
  \begin{cases}
    a)\ \mathfrak{F} \hookrightarrow \mathfrak{P}(\mathcal{M}) & F \mapsto \widetilde{F} = \mathrm{Dépl}(F) \subset \mathcal{M} \\
    b)\ \mathfrak{F} \hookrightarrow \mathfrak{P}(\mathfrak{P}(\mathcal{M})) & \text{plus précisément } \mathfrak{F} \to \mathrm{Fig}(\mathcal{M}), \\
     & F \mapsto \mathrm{Fig}_{\mathcal{M}}(F) \subset \mathfrak{P}(\mathcal{M})
  \end{cases}
\]
\marginal{b) se déduit de a) et de 1.2. a)}
\marginal{\ill{} \ill{}}
\note{la première note marginale est écrite en oblique dans la marge gauche ; la seconde, de deux mots, est sous b), une flèche la rattache à « b) » et une autre flèche courbe la mène à « $F \mapsto \mathrm{Fig}_{\mathcal{M}}(F)$ » ; dans cette dernière formule le « $\mathrm{F}$ » de « $\mathrm{Fig}$ » est écrit en surcharge sur un $\widetilde{F}$}

où $\mathrm{Fig}(\mathcal{M})$ désigne la partie de $\mathfrak{P}(\mathfrak{P}(\mathcal{M}))$ formée des figures ensemblistes dans $\mathcal{M}$. \struck{On} On peut considérer que la première application correspond à une relation entre $\mathcal{M}$ et $\mathfrak{F}$, appelée relation d'incidence. Pour une figure $F$, $\widetilde{F}$ s'appelle l'ensemble des \emph{multistrates} \emph{incidentes}, ou le \emph{déploiement} de la figure $F$. Si $X \in \mathcal{M}$, $F \in \mathfrak{F}$, on dira que la multistrate $X$ est \emph{incidente} à la figure $F$, \struck{de la figure $F$} ou encore que c'est une \emph{strate} de $F$, si $X \in \widetilde{F}$. D'autre part, tout\struck{e} \add{élément $X$} \struck{de} de $\mathcal{M}$ (i.e.\ \struck{toute} multistrate), regardé\supplied{e} comme une figure par (1.2), a donc un déploiement $\widetilde{X}$, et \struck{la multistrate} \struck{mais} on pose
\[
  (1.4) \qquad \mathrm{Fig}_{\mathcal{M}}(F) = \lbrace \widetilde{X} \mid X \in \widetilde{F} = \mathrm{Dépl}(F) \rbrace \subset \mathfrak{P}(\mathcal{M})
  \quad \text{i.e.} \quad \mathrm{Fig}_{\mathcal{M}}(F) \in \mathfrak{P}(\mathfrak{P}(\mathcal{M}))
\]
et il résultera des axiomes que c'est une figure ensembliste dans $\mathcal{M}$, indexée fidèlement par l'ens.\ $\widetilde{F}$ des strates de $F$. \struck{\ill{}} En fait, $\mathcal{M}$ sera muni d'une relation d'ordre $\leq$, cf.\ plus bas, et $\widetilde{F}$ \add{$\subset \mathcal{M}$} sera une partie fermée de $\mathcal{M}$, et pour tout $X \in \widetilde{F}$, \struck{$\widetilde{X}$} on aura
\[
  (1.5) \qquad \widetilde{X} = \lbrace Y \in \mathcal{M} \mid Y \leq X \rbrace .
\]

\page{3}
\note{son numéro « 3 » en haut de la page}
À cause de cette interprétation, le passage de $\widetilde{F} \subset \mathcal{M}$ à $\mathrm{Fig}_{\mathcal{M}}(F)$ est à tel point évident, que \add{la considération de} cette figure ensembliste dans $\mathcal{M}$ ne semble \uncertain{nullement} importante — mais à voir….

\struck{Quand on identifie}
On pourra utiliser les deux applications (1.3) pour identifier une figure, suivant les cas, à une partie de $\mathcal{M}$ [via $F \mapsto \widetilde{F}$] ou à une figure ensembliste dans $\mathcal{M}$ [via $F \mapsto \mathrm{Dépl}(F)$]. Dans le premier cas, \uncertain{compte tenu de} $\mathcal{M} \hookrightarrow \mathfrak{F}$ (1.2), par la formule 1.5. Cette dernière inclusion s'identifie à
\[
  X \longmapsto \widetilde{X} = \lbrace Y \in \mathcal{M} \mid Y \leq X \rbrace \qquad \mathcal{M} \to \mathfrak{P}_{\mathrm{ferm}}(\mathcal{M}) .
\]
\note{dans le second crochet, « $\mathrm{Dépl}(F)$ » est sur la page ; d'après (1.3), $\mathrm{Dépl}(F) = \widetilde{F}$, et c'est $\mathrm{Fig}_{\mathcal{M}}(F)$ qu'on attend}
Donc il faudra être prudent \ill{} l'identification dont il a été question, des éléments $X \in \mathcal{M}$ multistrates à des figures particulières, quand celles-ci sont interprétées comme des parties de $\mathcal{M}$ : la figure associée à $X \in \mathcal{M}$ n'est nullement $\lbrace X \rbrace$, mais $\widetilde{X}$ = ens.\ des multistrates majorées par $X$ pour $\leq$ (ou encore, des sous-multistrates, ou multistrates incidentes à $X$, ou encore des \struck{\ill{}} \emph{strates} de $X$).

\page{4}
\note{son numéro « 4 » en haut de la page}
D'autre part, il y a deux \add{autres} applications
\[
  (1.6) \qquad
  \begin{cases}
    a)\ \mathfrak{F} \to \mathfrak{P}(\mathcal{M}) & F \mapsto \mathrm{Omb}_{\mathcal{M}}(F) \\
    b)\ \mathfrak{F} \hookrightarrow \mathfrak{P}\mathfrak{P}(\mathcal{M}) \ \text{et} = \mathfrak{F} \to \mathrm{Fig}(\mathcal{M}) \subset \mathfrak{P}\mathfrak{P}(\mathcal{M}) & F \mapsto \mathrm{Multomb}_{\mathcal{M}}(F) ,
  \end{cases}
\]
\note{« (1.6) » est écrit en surcharge sur « (1.4) »}
où $\mathrm{Omb}_{\mathcal{M}}(F)$ s'appelle l'\emph{ombre de la figure} $F$ sur $\mathcal{M}$, ou la \emph{grande ombre} \emph{de} $F$. Une multistrate \add{$X$} qui est dans la grande ombre d'une figure $F$ est dite « \emph{raffiner} » \emph{la figure} $F$ (c'est cette fois déduit d'une relation d'ordre $\ll$ sur $\mathfrak{F}$, la raffinement, par la condition $X \ll F$). Il n'est pas clair que cette application $\mathfrak{F} \to \mathfrak{P}(\mathcal{M})$ soit injective, i.e.\ que $F$ s'identifie à un ens.\ de parties de $\mathcal{M}$, via cette application. Par contre il en sera ainsi de (1.6) b), qui est déduit de a) via \struck{la} formation de $\widetilde{F}$, comme (1.3 b) est déduit de (1.3.\ a) :
\[
  (1.7) \qquad \mathrm{Multomb}(F) = \lbrace X \in \widetilde{F} \mid \mathrm{Omb}\, X \rbrace ,
\]
\note{« (1.7) » remplace un « (1.5) » biffé ; la formule est ainsi sur la page, on attend $\lbrace \mathrm{Omb}\, X \mid X \in \widetilde{F} \rbrace$, par analogie avec (1.4)}
alors que $\mathrm{Omb}(X)$ s'explicite par
\[
  (1.8) \qquad \mathrm{Multomb}(X) = \lbrace Y \in \mathcal{M} \mid Y \ll X \rbrace ,
\]
\note{à gauche, « $\mathrm{Mult}$ » précède un « $\mathrm{O}$ » repassé ; le texte annonce $\mathrm{Omb}(X)$}
analogue à (1.5), où cette fois on a utilisé la relation \add{d'ordre} $\ll$ au lieu de $\leq$. \struck{C'est la}

\page{5}
\note{son numéro « 5 » en haut de la page}
\struck{une figure de}
Pour toute figure $F$, $\mathrm{Multomb}(F)$ \add{est appelé} \struck{dans} la \emph{grande} \emph{multiombre} de $F$, ou la $\mathcal{M}$-multiombre de $F$. \struck{On} C'est en fait une figure ensembliste dans $\mathcal{M}$, \add{indexée fidèlement par $\widetilde{F}$ (le déploiement de $F$)} et dont la connaissance implique celle de $F$. \add{(Jusqu'ici \uncertain{puisque} il n'était pas question de $\mathcal{L}$.)} Utilisant \struck{d'autre part} \add{maintenant} l'application
\[
  \mathcal{L} \hookrightarrow \mathcal{M} ,
\]
et les applications images inverses correspondantes
\[
  \mathfrak{P}(\mathcal{M}) \to \mathfrak{P}(\mathcal{L}) , \qquad \mathrm{Fig}(\mathcal{M}) \to \mathrm{Fig}(\mathcal{L}) ,
\]
on trouve, en composant avec 1.6 a) b)
\[
  (1.9) \qquad
  \begin{cases}
    a)\ \mathfrak{F} \to \mathfrak{P}(\mathcal{L}) & F \mapsto \mathrm{omb}(F) \\
    b)\ \mathfrak{F} \to \mathrm{Fig}(\mathcal{L}) & F \mapsto \mathrm{multomb}(F)
  \end{cases} ,
\]
\note{les deux flèches de (1.9) sont écrites en surcharge sur des flèches à crochet}
dont la première n'a aucune raison d'être injective, le second par contre a une tendance à l'être. \struck{On va} \add{On} appelle $\mathrm{omb}(F)$ la \emph{petite ombre} de $F$, ou simplement son ombre quand une confusion semble exclue, et \struck{de même} \uncertain{de} $\mathrm{multomb}$, la (\emph{petite}) \emph{multiombre de la} figure $F$.

\page{6}
\note{son numéro « 6 » en haut de la page}
Ainsi, l'ombre d'une figure est formée des lieux qui la raffinent — on dit aussi que ce sont les \emph{lieux « de la figure »} (\ill{} un peu \uncertain{abusivement}, ou les strates de la figure qui sont des lieux). La multiombre est donc formée des ombres de ses strates, elle est indexée par le déploiement $\widetilde{F}$, ensemble des strates de $F$, la relation d'\struck{\ill{}}\add{inclusion} de la figure ensembliste n'étant autre que la relation d'incidence entre strates de $F$, i.e.\ la relation dans $\widetilde{F}$ induite par $\leq$ dans $\mathcal{M}$ (ou par $\ll$ au choix, se réduisant au $=$). Alors qu'une figure n'est reconstituable pratiquement jamais à partir de son ombre, elle est reconstituable pourtant \struck{donc} \uncertain{dans} (cas des « Ateliers ensemblistes ») \struck{à} à partir de sa multiombre. Nous y reviendrons.

\struck{§} Il y a lieu de parler des relations d'ordre canoniques

\page{7}
\note{son numéro « 7 » en haut de la page}
\[
  (1.10) \qquad \leq,\ \ll,\ \preccurlyeq \quad \text{relations d'ordre dans } \mathfrak{F}
\]
et les relations induites dans $\mathcal{M}$
\[
  (1.11) \qquad \leq,\ \ll \quad \text{dans } \mathcal{M}
\]
[et la relation d'ordre $\preccurlyeq$ induite dans $\mathcal{M}$ est moins importante, comme structure \uncertain{spécifique} \ill{} $\mathcal{M}$].
\note{le troisième signe, que l'on rend par $\preccurlyeq$, est un $\leq$ dont le trait inférieur est courbé}

\marginal{\emph{NB} Les relations d'ordre induites par $\leq$, $\ll$ sur $\mathcal{L}$ sont discrètes.}
\marginal{\emph{NB} $G \leq F \Rightarrow G \ll F$ : une sous-figure de $F$ raffine $F$}
\note{ces deux notes sont écrites en oblique dans la marge gauche, en regard de (1.11) et des définitions qui suivent}

$G \leq F$ \quad $G$ est dit une \emph{sous-figure} de $F$ (si \struck{$G$} \add{$G = X$} $\in \mathcal{M}$, c'est donc une \emph{strate} de $F$)

$G \ll F$ \quad $G$ est dit un \emph{raffinement} de $F$ (si $G = X \in \mathcal{M}$, c'est donc un élément de la grande ombre de $F$)

$G \preccurlyeq F$ \quad $G$ est dit une \emph{subdivision} de $F$, on dit que la figure $G$ subdivise la figure $F$.

Je note enfin \struck{deux} autres relations importantes \struck{dans} entre figures :
\[
  (1.12) \qquad
  \begin{cases}
    F \text{ et } G \ \underline{\text{compatibles}} \\
    F \text{ et } G \ \underline{\text{disjoints}}
  \end{cases} ,
\]
qui l'une et l'autre peuvent s'exprimer en termes de la relation d'ordre $\leq$ :
$F$ et $G$ compatibles : ce sont des sous-figures d'une même figure, i.e.\ $\lbrace F, G \rbrace$ majoré pour $\leq$ ;
$F$ et $G$ disjoints : $F$ et $G$ compatibles, et leur

\page{8}
\note{son numéro « 8 » en haut de la page}
inf dans $\mathfrak{F}$ est $\emptyset_{\mathfrak{F}}$, le plus petit élément de $\mathfrak{F}$ ou \emph{figure vide}. Ainsi
\[
  F \text{ et } G \text{ disjoints} \Longrightarrow F \text{ et } G \text{ compatibles.}
\]
On a, par les axiomes :
\[
  (1.13) \qquad F \text{ et } G \text{ disjoints} \iff
\]
\struck{tout} $\mathrm{omb}(F)$ et $\mathrm{omb}(G)$ \struck{$\mathrm{omb}(F) \cap \mathrm{omb}(G) = \emptyset$} sont « disjoints » en un sens suivant : ils sont disjoints ensemblistement i.e.\ $\mathrm{omb}(F) \cap \mathrm{omb}(G) = \emptyset$, et tout élément de $\mathrm{omb}(F)$ est compatible avec tout élt de $\mathrm{omb}(G)$)
\note{ce membre de droite de (1.13) est écrit en colonne à droite de la page}

— i.e.\ deux figures sont disjointes ssi elles n'ont pas de lieu commun, et tout lieu de l'une est \struck{disjoint} \add{\struck{co} disjoint} (au sens de \struck{la théorie} l'atelier) de \struck{tout} lieux de l'autre ;

Cela montre que les relations de \struck{dis}compatibilité et de disjonction sont importantes aussi dans les \struck{\ill{}} ensembles $\mathcal{M}$, $\mathcal{L}$ de $\mathfrak{F}$. Dans le cas des ateliers ensemblistes, ces relations sont « triviales » sur $\mathcal{L}$ : deux lieux sont \struck{disjoints} \add{compatibles}, et disjoints ssi ils sont distincts. Par contre, même dans le cas ensembliste et « modéré », la relation

\page{9}
\note{son numéro « 9 » en haut de la page}
de \emph{compatibilité} dans $\mathcal{M}$ est une \struck{relation} relation importante, qui se surajoute aux relations $\leq_{\mathcal{M}}$, $\ll_{\mathcal{M}}$. Moralement, pour récupérer l'atelier (à iso unique près) il suffit de connaître sur $\mathcal{M}$ ces deux relations d'ordre, et la relation de compatibilité.
\marginal{au moins \ill{} dans le cas \uncertain{modéré}}
\note{note écrite en oblique dans la marge gauche, en regard de ces lignes}

Les relations \add{(compatibilité, disjonction)} \struck{\ill{}} sont si importantes qu'elles méritent chacune une notation lapidaire. Je propose \struck{\ill{}}
\[
  (1.14) \qquad F \parallel G
\]
pour la disjonction, et \struck{\ill{}}
\[
  (1.15) \qquad F \mathbin{\Diamond} G \quad \text{ou} \quad F \mathbin{\between} G
\]
pour la compatibilité (NB $\between$ est plus rapide à écrire que $\Diamond$, $\Leftrightarrow$, mais est moins suggestif de l'idée combinatoire exprimée.
\note{les deux notations biffées devant « $F \parallel G$ » et « $F \mathbin{\Diamond} G$ » sont surchargées et raturées. Le signe rendu $\Diamond$ est un losange partagé par un trait vertical ; le signe rendu $\between$, comme au lot 2, est un sablier (deux chevrons opposés par le sommet, fermés en haut et en bas), et « $F \between G$ » est entouré d'un cercle ; le signe rendu $\Leftrightarrow$ est une double flèche horizontale inscrite dans un losange}

On n'a rien dit pour l'instant au sujet de la relation \struck{$\ll$} $\preccurlyeq$ de \uncertain{raffinement} sur $\mathfrak{F}$. Elle est formellement la relation la plus délicate, et est la relation la plus spécifique dans l'« Algèbre des figures ». Sa signification géométrique s'éclaire par les exemples, et les axiomes qui aux
\note{à la page 7, $\preccurlyeq$ est la \emph{subdivision} et $\ll$ le \emph{raffinement} ; le mot lu ici « raffinement » est resté après que $\ll$ a été biffé et remplacé par $\preccurlyeq$}

\page{10}
\note{son numéro « 10 » en haut de la page}
ci inspirent.

\emph{Axiomes} On peut prendre pour ens.\ de base d'une axiomatique des ateliers, soit $\mathfrak{F}$, soit $\mathcal{M}$, soit $\mathcal{L}$. On trouve des \ill{} les plus simples sur $\mathfrak{F}$, les \struck{plus} sophistiqués sur $\mathcal{L}$. Par contre, les axiomes \uncertain{naturels} sont les moins transparents quand on reste plus ou moins dans $\mathfrak{F}$, et les plus transparents quand on travaille dans $\mathcal{L}$ (ou peut-être pareil pour $\mathcal{M}$ ?).

\emph{Sur} $\mathfrak{F}$, on peut prendre comme structure de base, au choix
\[
  (1.16) \qquad
  \begin{cases}
    a)\ \leq,\ \ll \\
    b)\ \leq,\ \preccurlyeq \\
    c)\ \ll,\ \text{et}\ \mathfrak{F} \to \mathrm{P}(\mathfrak{F})
  \end{cases}
\]
\marginal{(20 juin) Non, \ill{} nouveau \ill{} \ill{} que pour $\ll$}
\note{ajout daté, de sa main, écrit à droite de b) et relié à b) par une flèche ; la date « 20 juin » est la sienne}
\note{dans c), un signe biffé avant « et »}
(notion $F \mapsto \widetilde{F}$, le déploiement de $F$ vu comme partie de $\mathfrak{F}$) \struck{i.e.\ la \struck{\ill{}} relation d'inc\supplied{idence}} ; sous c), i.e.\ la \emph{relation d'incidence} « $F$ est une strate de $G$ » (i.e.\ $F \in \mathcal{M}$ \emph{et} $F \leq G$).
\note{le commentaire entre parenthèses est écrit à droite de c) ; « i.e.\ la relation d'incidence… » est écrit sous c) et rattaché par une accolade à « $\ll$, et $\mathfrak{F} \to \mathrm{P}(\mathfrak{F})$ »}

\struck{Je} Avant de préciser ce point, je vais grouper dans un diagramme les six relations fondamentales dans $\mathfrak{F}$, en \uncertain{incluant}
\[
  (1.17) \qquad F \mathrel{\triangleleft} G \overset{\mathrm{dfn}}{\iff} \text{« $F$ incident à $G$ » i.e. } F \leq G \text{ et } F \in \mathcal{M}
\]

\page{11}
\note{son numéro « 11 » en haut de la page}
(1.18)

\begin{tikzcd}[column sep=small]
 & & G \mathrel{\triangleleft} F \arrow[d, Rightarrow] & & \\
G \preccurlyeq F \arrow[dr, Rightarrow] & & G \leq F \arrow[dl, Rightarrow] \arrow[dr, Rightarrow] & & G \parallel F \arrow[dl, Rightarrow] \\
 & G \ll F & & G \between F &
\end{tikzcd}

\note{les flèches sont doubles (implications) ; un cadre oblique, tracé à la main, entoure $G \leq F$ et $G \ll F$ avec la flèche qui les joint}

Parmi les six relations, il y a quatre relations d'ordre $\triangleleft$, $\leq$, $\preccurlyeq$, $\ll$, et les deux autres relations $\parallel$ et $\between$ sont symétriques. \struck{Les} La relation $\between$ est \add{de plus,} réflexive, mais $\parallel$ ne l'est que si $\mathfrak{F}$ est réduit à son seul élément $\emptyset_{\mathfrak{F}}$, car
\[
  (1.19) \qquad F \parallel F \iff F = \emptyset_{\mathfrak{F}} \quad \text{figure vide.}
\]

\marginal{Non, $\triangleleft$ n'est pas rel.\ d'ordre, elle est transitive mais non réflexive, $F \mathrel{\triangleleft} F \Leftrightarrow F \in \mathcal{M}$. Pour \ill{} une rel.\ d'ordre \ill{} avoir écrit $F \trianglelefteq F \overset{\text{déf}}{\Longleftrightarrow} F \mathrel{\triangleleft} F$ ou $F = F$}
\note{note écrite en oblique dans la marge gauche, en regard de l'énumération des « quatre relations d'ordre » ; le dernier $F \trianglelefteq F$ (triangle souligné) est aussi porté seul au bas de la note. La définition proposée, avec « $F = F$ », est telle quelle sur la page ; on attend $F \mathrel{\triangleleft} G$ ou $F = G$}

\emph{Remarques} a) les relations $\triangleleft$ et $\leq$ se déterminent mutuellement. Connaissant $\leq$, on définit $\mathcal{M} \subset \mathfrak{F}$ comme
\[
  (1.20) \qquad \mathcal{M} = \lbrace X \in \mathfrak{F} \mid X \text{ strict.\ irréductible pour } \leq ,\ \text{i.e. } X \leq \textstyle\sup_i F_i \Rightarrow \exists\, i \text{ avec } X \leq F_i \rbrace
\]
puis $\triangleleft$ par (1.17). \struck{D'autre} D'autre part, connaissant $\triangleleft$ on définit $\mathcal{M}$ par
\[
  (1.21) \qquad \mathcal{M} = \lbrace X \in \mathfrak{F} \mid X \mathrel{\triangleleft} X \rbrace
\]
puis, pour $F \in \mathfrak{F}$, son déploiement
\[
  (1.22) \qquad \widetilde{F} = \lbrace X \in \mathcal{M} \mid X \mathrel{\triangleleft} F \rbrace
\]
et enfin
\[
  (1.23) \qquad F \leq G \iff \widetilde{F} \subset \widetilde{G}
\]
\note{dans (1.20), la dernière inégalité, « $X \leq F_i$ », est surchargée ; dans (1.22), la page écrit « $X \mathrel{\triangleleft} \mathfrak{F}$ », pour $X \mathrel{\triangleleft} F$}
Donc, pour déterminer \struck{la structure, il est équivalent} \add{les données \uncertain{minimales} qui} fixent la structure d'atelier, il est indifférent de s'y

\page{12}
\note{son numéro « 12 » en haut de la page}
prendre soit $\leq$, soit $\triangleleft$.

\struck{b) \ill{} et $\ll$ déterminent \ill{} $\preccurlyeq$, car}

\struck{$G \preccurlyeq F \iff G \ll F$}
\note{deux lignes biffées, la seconde surchargée}

b) Chacune des relations $\leq$, $\ll$ détermine la relation $\parallel$. En effet
\[
  (1.24) \qquad
  \begin{cases}
    a)\ F \parallel G \iff \lbrace F, G \rbrace \text{ majoré pour } \leq, \text{ et } \mathrm{Inf}_{\leq}(F,G) = \emptyset_{\mathfrak{F}} \\
    b)\ F \parallel G \iff \mathrm{Sup}_{\ll}(F,G) \text{ existe, et } \mathrm{Inf}_{\ll}\, F,G = \emptyset_{\mathfrak{F}}
  \end{cases}
\]
\marginal{Vérifier $\rightarrow$}
\note{la note, dans la marge gauche, pointe vers b) ; dans a), le premier $\iff$ est écrit en surcharge sur un autre signe}
(NB $\emptyset_{\mathfrak{F}}$ est à la fois plus petit élément de $\mathfrak{F}$ pour $\leq$, et pour $\ll$)

c) La relation $\ll$ détermine $\preccurlyeq$, par
\[
  (1.25) \qquad F \preccurlyeq G \iff F \ll G, \text{ et } \forall H \in \mathfrak{F} \text{ on a } (H \parallel F \iff H \parallel G)
\]
\note{une croix dans la marge gauche, en regard de (1.25)}
et on applique b).

d) La relation $\ll$ détermine $\mathcal{L} \subset \mathfrak{F}$ comme l'ens.\ des éléments minimaux de $\mathfrak{F} \smallsetminus \lbrace \emptyset_{\mathfrak{F}} \rbrace$ :
\[
  (1.26) \qquad \mathcal{L} = \lbrace X \in \mathfrak{F},\ X \neq \emptyset_{\mathfrak{F}} \mid Y \in \mathfrak{F},\ Y \ll X,\ Y \neq \emptyset_{\mathfrak{F}} \Rightarrow Y = X \rbrace
\]
(Chaque lieu est une \emph{figure} non vide dont le seul raffinement non vide est elle-même).

\struck{On dit $\ll$ implique la relation $\between$ par}
\struck{(1.27) $F \between G \iff \exists\, V$}
\note{la ligne est biffée, et la formule, encadrée, est annulée par des hachures obliques}
Je ne vois pas que $\ll$ (+ $\preccurlyeq$, c'est kif-kif) détermine $\leq$, ni même $\between$. Mais

(e) $\ll$ et $\between$ déterminent $\leq$ par
\[
  (1.27) \qquad G \leq F \iff (G \ll F \text{ et } G \between F)
\]
\note{« e » est cerclé ; une croix dans la marge gauche, en regard de (1.27)}
Donc finalement,

\page{13}
\note{son numéro « 13 » en haut de la page}
on trouve trois possibilités optimales de définir la structure d'atelier sur $\mathfrak{F}$ :
\[
  (1.28) \qquad
  \begin{cases}
    \text{soit par } \boxed{\leq,\ \ll} \\
    \text{soit par } \leq,\ \preccurlyeq \\
    \text{soit par } \ll,\ \between .
  \end{cases}
\]
\marginal{Non, ça ne suffit pas, d'après nouveau point de vue…}
\note{ajout écrit à droite de la deuxième ligne de (1.28) ; cf.\ l'ajout daté du 20 juin à la page 10}

Dans la première mouture, j'ai choisi $\leq$, $\preccurlyeq$, en définissant en termes de ceux-ci $\ll$ \add{est} la relation de \uncertain{préordre} engendrée par $\leq$, $\preccurlyeq$, donc on a
\[
  (1.29) \qquad G \ll F \iff \exists\, F' \text{ avec } G \leq F' \preccurlyeq F
\]
[correspondant au \uncertain{contexte} de \ill{} \add{\ill{}} \struck{\ill{}} intuitif], qui est bien une notion de subdivision, plutôt \ill{} celle de raffinement. Mais techniquement c'était un peu laborieux, de déduire les propriétés de transitivité et \add{surtout} de réflexivité pour $\ll$, et ça amenait à multiplier les axiomes en nombre. Toujours techniquement, il semblerait que ce soit en partant de $\leq$, $\ll$ que soit le plus simple. C'est l'approche que je vais suivre maintenant. Quant à la possibilité de passer par $\ll$, $\between$, c'est plutôt une curiosité qu'autre chose.
\note{la « première mouture » est le chapitre IV, dossier 156-4. La page écrit « relation de … engendrée par $\leq$, $\preccurlyeq$ » ; deux ordres engendrent un préordre, d'où la lecture proposée}

\marginal{J'ai \ill{} le \ill{} de (1.29) \ill{} le 20, \ill{} plus que \ill{} \ill{} $\ll$ \ill{} $\leq$, $\preccurlyeq$ \ill{} $\ll$ \ill{} et \ill{}}
\note{longue note écrite en oblique dans la marge gauche, en regard de (1.29) et du paragraphe qui suit, en grande partie illisible ; « le 20 » y renvoie sans doute au 20 juin, date de l'ajout de la page 10}

\page{14}
\note{son numéro « 14 » en haut de la page}
En résumé, en termes de $\mathfrak{F}$, c'est la donnée de deux \emph{relations} \emph{d'ordre}
\[
  (1.30) \qquad \boxed{\leq,\ \ll \ \text{(rel.\ d'ordre dans } \mathfrak{F}) \ \text{telles que} \ F \leq G \Rightarrow F \ll G}
\]
que je vais partir.
\note{la formule est encadrée ; « rel.\ d'ordre dans $\mathfrak{F}$ » est écrit sous « $\leq$, $\ll$ », rattaché par une accolade}

En termes de celles-ci, on définit \struck{$\mathcal{M}$}
\[
  \mathcal{M} \subset \mathfrak{F}
\]
par (1.20), \uncertain{avec} les relations $\leq_{\mathcal{M}}$, $\ll_{\mathcal{M}}$ induites par $\leq_{\mathfrak{F}}$, $\ll_{\mathfrak{F}}$, puis
\[
  \mathcal{L} \subset \mathcal{M}
\]
par (1.25). L'application
\[
  (*) \qquad \mathfrak{F} \hookrightarrow \mathfrak{P}(\mathcal{M}) \qquad F \mapsto \widetilde{F} = \mathrm{Dépl}(F)
\]
est \add{donc} définie par
\[
  (1.31) \qquad \widetilde{F} = \lbrace X \in \mathcal{M} \mid X \leq F \rbrace
\]
et
\[
  (**) \qquad \mathfrak{F} \hookrightarrow \mathrm{Fig}(\mathcal{M}) \qquad F \mapsto \mathrm{Fig}_{\mathcal{M}}(F)
\]
par (1.4) \struck{(1.32)}
\[
  \mathrm{Fig}_{\mathcal{M}}(F) = \lbrace \widetilde{X} \mid X \in \widetilde{F} \rbrace \subset \mathfrak{P}(\mathcal{M})
  \quad \text{avec} \quad \widetilde{X} = \lbrace Y \in \mathcal{M} \mid Y \leq X \rbrace .
\]
\note{« (1.25) » est sur la page ; $\mathcal{L}$ est défini en (1.26), page 12. Les marques « (*) » et « (**) » sont écrites en surcharge ; dans (*), un $X$ biffé précède $F \mapsto \widetilde{F}$}

Ainsi, $\mathfrak{F}$ est interprété comme un ens.\ de parties dans $\mathcal{M}$ via (*), et $\leq$ s'interprète simplement par l'inclusion
\[
  (1.32) \qquad F \leq G \iff \widetilde{F} \subset \widetilde{G}
\]

\page{15}
\note{son numéro « 15 » en haut de la page}
Par contre la relation $\ll$ n'a pas d'interprétation simple évidente, ni en termes des \add{déploiements} $\widetilde{F} \subset \mathcal{M}$, ni en termes des multidéploiements $\mathrm{Fig}(F)$ ; ainsi $G \ll F$ n'implique nullement $\mathrm{Fig}_{\mathcal{M}}(G) \ll \mathrm{Fig}_{\mathcal{M}}(F)$. Aussi l'application (**) et la notation $= \mathrm{Fig}_{\mathcal{M}}(F)$ sont-ils finalement un peu bidon !

En termes de $\mathcal{M}$, la structure est donnée par
\[
  (1.33) \qquad
  \boxed{
  \begin{array}{l}
    \mathfrak{F} \subset \mathfrak{P}(\mathcal{M}) \ \text{ens.\ de parties de } \mathcal{M} \\
    \ll \ \text{relation d'ordre dans } \mathfrak{F}, \text{ impliquée par l'inclusion (notée } F \leq G)
  \end{array}}
\]
Ça a l'air moins sympathique du point de vue simplicité des données et d'autant plus qu'il n'y a pas \add{à première vue} \add{d'interprétation directe} de $\ll$ en termes d'\struck{\ill{}} \add{ensemblistes} \ill{}. Mais c'est ici que vont servir les (grandes) ombres ! (et avant d'introduire $\ll$)

En termes de la \add{première} donnée (1.33), le \struck{\ill{}} premier axiome est \struck{\ill{}} le fait que $\mathfrak{F}$ est \struck{une famille d'} \struck{une figure} ens.\ \add{\emph{préadmissible}} \struck{dans} $\mathcal{M}$ de \struck{Axiome} \uncertain{parties de} $\mathcal{M}$, \uncertain{donc} dans la donnée de $\mathfrak{F}$ \add{implique} \struck{revient à la donnée} d'une relation de préordre dans $\mathcal{M}$ (la rel.\ $\leq$) — lequel se
\marginal{attention au « pré » !}
\note{la note marginale est en regard de « préadmissible », dont le « pré » est souligné deux fois ; au-dessous, dans la marge, une autre note est annulée par des hachures. Les ratures et les ajouts de ces dernières lignes s'enchevêtrent et leur ordre n'est pas sûr. Un « 16 » est écrit sous l'accolade en fin de page}

\page{16}
\note{son numéro « 16 » en haut de la page}
une relation d'ordre ; en d'autres termes

(i) $\forall X \in \mathcal{M}$, dans l'ens.\ $\mathfrak{F}^{X}$ des $F \in \mathfrak{F}$ telles que $X \in F$, il y a un plus petit élément, soit $F_X$ (NB ce sera l'ens.\ des $Y \in \mathcal{M}$ tels que $Y \leq X$)

(ii) L'application $X \mapsto F_X$ est injective, i.e.\ si $X, Y \in \mathcal{M}$, $X \neq Y$, alors $\exists\, F \in \mathfrak{F}$ telle que $X \in F$, $Y \notin F$, ou l'inverse $Y \in F$, $X \notin F$.

Dès lors, les $F \in \mathfrak{F}$ \struck{\ill{}} \add{\uncertain{sont les} \ill{}} contenant les $\mathcal{M}_{\leq X}$, $X \in \mathcal{M}$ \struck{\ill{}} \add{\ill{}} \add{fermées} de \uncertain{plus} \ill{} suppose de $\mathcal{M}$, \struck{\ill{} donnée \ill{} $\mathfrak{F} \subset \mathfrak{P}(\mathcal{M})$ satisfaisant (i) (ii), revient \ill{} que \ill{}} $\emptyset$ (partie vide de $\mathcal{M}$) $\in \mathfrak{F}$ (ce qui implique que $\mathfrak{F}$ n'est \emph{jamais} admissible !).
\note{passage très raturé : deux lignes sont biffées et plusieurs ajouts interlinéaires s'y enchevêtrent ; on donne ce qui se lit}

Dans GF IV (Analysis situs, première mouture) j'ai explicité en long et en large la relation entre la donnée de $(\mathfrak{F}, \leq)$, satisfaisant les conditions qui y étaient \uncertain{appelées} \emph{C 1}, \emph{C 2}, \emph{C 3} (\uncertain{voir} §§~5, 6, 8), et celle d'un ens.\ ordonné $\mathcal{M}$, muni d'un ens.\ de parties fermées $\mathfrak{F}$ \add{qui} contenant les $\mathcal{M}_{\leq X}$, la partie vide, formé de parties fermées (pour l'ordre de $\mathcal{M}$), plus la condition : si $F, G \in \mathfrak{F}$, tels que $\forall X \in F$, $Y \in G$ on ait $\mathcal{M}_{\leq X} \cup \mathcal{M}_{\leq Y} \in \mathfrak{F}$, alors $F \cup G \in \mathfrak{F}$. (Condition de compatibilité).
\marginal{J'appellerai ces trois conditions plutôt \emph{At 1}, \emph{At 2}, \emph{At 3}}
\note{note écrite en oblique dans la marge gauche, en regard de « C 1, C 2, C 3 ». Le renvoi est au chapitre IV, dossier 156-4 ; les numéros de paragraphes, surchargés, se lisent « 5, 6, 8 » sans certitude}

\page{17}
\note{son numéro « 17 » en haut de la page}
Il n'y a pas à y revenir. Je vais plutôt m'attacher à présent aux interprétations en termes d'ombres (grandes ou petites). Il est entendu que dans un premier temps, je m'attache aux aspects « algébriques » ou « triviaux » d'un formalisme des figures, et un contexte de \uncertain{dignes} de bon raisonnement, de telle façon que

1°) Ceux-ci \struck{\ill{}} valables dans tout « atelier » de la forme $\mathrm{Fig}(\text{\struck{\ill{}}}\, L)$, formé des figures ensemblistes dans un ensemble donné $L$ (sans conditions de finitude d'aucune sorte), et

2°) qu'ils soient applicables aussi, disons, aux figures \emph{simpliciales} affines finies, dans un espace affine sur $\mathbb{R}$ (ou, plus généralement, sur un corps ordonné quelconque).

Pour expliciter, en termes de l'ens.\ ordonné $(\mathcal{M}, \leq)$, les ens.\ $\mathfrak{F}$ de figures, en tant qu'ens.\ de parties de $\mathcal{M}$, \struck{il est} on dira \emph{figures de type fini}
\[
  (1.34) \qquad F \in \mathfrak{F} \text{ de t.f.} \iff \widetilde{F} = \text{réunion finie d'ens.\ } \mathcal{M}_{\leq X_i}
\]
(pour un ens.\ fini de strates de \struck{$\mathcal{M}$})
\note{la parenthèse est écrite sous la formule ; dans (1.34), un mot biffé précède « réunion »}

\page{18}
\note{son numéro « 18 » en haut de la page}
la chose essentielle est la connaissance de
\[
  (1.35) \qquad \mathbin{\between_{\mathcal{M}}} \ \text{relation binaire sym.\ réfl.\ (sur } \mathcal{M})
\]
\struck{qui implique} satisfaisant
\[
  (1.36) \qquad X \between Y,\quad X' \leq X,\ Y' \leq Y \Longrightarrow X' \between Y'
\]
\struck{Alors} il y a corr.\ biunivoque entre l'ens.\ de telles relations sur $\mathcal{M}$, et d'un ens.\ $\mathfrak{F}$ \add{$\subset \mathfrak{P}(\mathcal{M})$} satisfaisant les conditions énumérées précédemment, plus \struck{celle disant} que les $F \in \mathfrak{F}$ soient \struck{\ill{}} de type fini. C'est en ce sens que j'ai pu dire que « moralement », le couple $(\mathcal{M}, \mathfrak{F} \subset \mathfrak{P}(\mathcal{M}))$ équivalait à \struck{\ill{}} la donnée de $\leq_{\mathcal{M}}$ et de $\between$ satisfaisant (1.36).
\note{cf.\ page 9 : « Moralement, pour récupérer l'atelier… il suffit de connaître sur $\mathcal{M}$ ces deux relations d'ordre, et la relation de compatibilité »}

\struck{Mais il s'agit maintenant}
Pratiquement, dans presque tous les exemples, c'est via $\mathcal{M}$ \add{(multistrates)} muni de $\leq$ et d'une relation de disjonction $\between$, qu'on définira les ateliers…
\note{« relation de disjonction $\between$ » est ainsi sur la page ; en (1.14)-(1.15), page 9, $\parallel$ note la disjonction et $\between$ la compatibilité, que désigne aussi (1.35)-(1.36)}

Il faut à présent en venir à la relation $\ll$. \struck{Comme elle} Elle ne pouvant l'instant s'exprimer surtout à formes des ombres, je commence par les grandes ombres

\page{19}
\note{son numéro « 19 » en haut de la page, écrit comme un « 17 »}
\[
  (*) \qquad \mathrm{Omb}(F) = \lbrace X \in \mathcal{M} \mid X \ll F \rbrace
\]
et la multiombre
\[
  (**) \qquad \mathrm{Multomb}(F) = \lbrace \mathrm{Omb}(X) \mid X \in \widetilde{F} \rbrace .
\]
\marginal{$\widetilde{F} = F$, si on identifie $F$ à une partie de $\mathcal{M}$ via $\mathfrak{F} \hookrightarrow \mathfrak{P}(\mathcal{M})$, $F \mapsto \widetilde{F}$}
\note{ajout à droite de (**), relié à $\widetilde{F}$ par un trait}

En apparence (*) utilise la connaissance de $\ll$ dans $\mathfrak{F}$ tout entier, mais il résulte des axiomes \struck{\ill{} qui dit} qu'elle ne dépend que de $\ll_{\mathcal{M}}$ (au lieu de $\ll_{\mathfrak{F}}$). En tous cas, il est clair \add{\ill{}} que la multiombre ne dépend \uncertain{que} de $\ll_{\mathcal{M}}$.
\marginal{NB Je vais \ill{} \ill{} \ill{} C 5 \ill{} C 4, \ill{} \ill{}}
\note{note écrite en oblique dans la marge gauche, en grande partie illisible}

Je veux que cette multiombre soit une \emph{figure ensembliste} dans $\mathcal{M}$, \add{et indexée fidèlement par $\widetilde{F}$}, \emph{appelée l'ombre}, (\ill{} \ill{} \ill{} \ill{} d'une \ill{} figure), \add{telle que $X \ll F$}, et
\marginal{avec la relation d'ordre $\leq$ dans $\mathfrak{F}$ correspondant à l'inclusion dans $\mathrm{Multomb}(F)$, i.e.\ dans $\mathfrak{F}$ ; C 4 si $X, Y \in \mathcal{M}$ tels que $X \ll Y$, $X \between Y$ \add{alors} $\Rightarrow X \leq Y$}
\note{note de la marge gauche, marquée d'un point cerclé qui la renvoie au texte ; « C 4 » y est souligné.}

\emph{C 5} si $X \in \mathcal{M}$, $F \in \mathfrak{F}$, (si $F^{X}$ est l'ens.\ des $Y \in \mathcal{M}$ tels que \struck{$X \leq Y$} $X \ll Y \leq F$ (l'ens.\ des strates de $F$ raffinées par $X$), on \struck{dans} (1.37) est que $F^{X}$ a un plus petit élément \add{pour $\leq$} (i.e.\ il existe une strate de $F$ raffinée par $X$, et parmi toutes les strates de $F$ raffinées par $X$, il en est une plus petite)
\note{la phrase est reprise en surcharge et les parenthèses ne se ferment pas toutes ; « (1.37) » est écrit dans la marge du texte de C 5}

\marginal{Cor Pour que si $X \in \mathcal{M}$, $F \in \mathfrak{F}$, $X \ll F$, il faut \ill{} il existe $Y \in \mathcal{M}$ \ill{} $X \ll Y \leq F$. Ainsi $X \ll F$ \ill{} \ill{} $Y \in F$ \ill{} $X \ll Y$}
\note{corollaire écrit en oblique dans la marge gauche, en regard de C 5 ; lecture partielle}

\struck{\emph{C 6} Soit \ill{} $X \in \mathcal{M}$, \ill{}. Alors $\exists\, Y \in \mathcal{M}$ tel que \add{i.e.\ $X \in F^{Y}$} a) $Y \ll X$ b) $Y \ll X' \leq X$, \add{$X' \in \mathcal{M}$} $\Rightarrow X' = X$ \struck{b) \ill{} $X$ est le plus petit élément de} (cf C 4)}
\note{« C 6 » et tout son énoncé sont biffés, la ligne de titre d'un trait, le reste de longs traits obliques}

\marginal{NB Le fait que dans $\mathrm{Multomb}(X)$, les $\mathrm{Omb}(X)^{\circ} \neq \emptyset$, est trivial, car $X \in \mathrm{Omb}(X)^{\circ}$}
\note{note au pied de la page, à gauche, sous la marge}

\page{20}
\note{son numéro « 20 » en haut de la page}
Ainsi on trouve donc une application
\[
  (*) \qquad \mathfrak{F} \longrightarrow \mathrm{Fig}(\mathcal{M}) , \qquad F \longmapsto \mathrm{Multomb}(F)
\]
Elle est injective, car \add{$X$ est connu quand $\mathrm{Omb}(X)$ est connu (c'est le plus grand élément de $\mathrm{Omb}(X)$ pour $\ll$), donc $\widetilde{F}$ est connu quand $\mathrm{Multomb}(F)$ l'est, donc $F$ aussi.} Il est immédiat
\[
  (1.38) \qquad G \leq F \iff \mathrm{Multomb}(G) \leq \mathrm{Multomb}(F)
\]
\note{l'explication de l'injectivité est ajoutée entre les lignes, à droite, et rattachée par une accolade ; le $\iff$ de (1.38) est écrit en surcharge. Sous les deux membres de droite, des accolades : « les strates sont les $\mathrm{Omb}(Y)$, $Y \in \widetilde{G}$ », « les strates sont les $\mathrm{Omb}(X)$, $X \in \widetilde{F}$ », et sous les deux : « or $\widetilde{G} \subset \widetilde{F}$ est une partie fermée de $\widetilde{F}$ si $G \leq F$ »}

\marginal{\emph{19 juin}}
\note{date de sa main, soulignée, en tête de l'alinéa}
Finalement, j'ai changé de point de vue pour le sens géométrique de la relation $G \ll F$, que je voyais comme « $G$ est une sous-figure d'une subdivision de $F$ ». Déjà dans \struck{\ill{}} le cas simplicial (i.e.\ la structure combinatoire des \add{$\mapsto$} multistrates étant \struck{celle} une structure de simplexe : l'ens.\ ordonné des sous-ens.\ d'un ens.\ fini…), cela n'a pas l'air naturel. Par contre, si on a un \uncertain{choix} modéré, et deux simplexes modérés $X$, $Y$ dans un $\mathbb{R}^n$ disons, on dira que $Y \ll X$ si c'est le cas pour les figures ensemblistes correspondantes — ce qui implique sans doute pas l'existence d'une subdiv.\ de $X$ telle que $Y$ en soit une multistrate de la
\note{dans « $Y \ll X$ », la première lettre est un $X$ surchargé, qu'on lit $Y$ d'après la suite. La date « 19 juin » précède ici celle du 20 juin portée page 10 : l'ajout de la page 10 est postérieur à la rédaction suivie. La page s'arrête au milieu de la phrase}

\end{document}
